\chapter{Forms, Conversion, Validation, and Problems}
\label{ch:forms-validation}
\chapterquote{A user must be allowed to type an invalid intermediate value without corrupting the last meaningful value.}{The form model rule}

\begin{chaptermap}
Core question & How does Corusco model editable forms without letting Swing fields own the business contract?\\
Primary APIs & \api{FormModel}, \api{FieldModel}, \api{TextFieldModel}, \api{ParseState}, \api{Problem}, \api{ProblemSet}, \api{ProblemTarget}, \api{ProblemSeverity}, \api{Validators}, \api{RuleSet}, and \api{AsyncFieldValidation}.\\
Plain Swing baseline & \api{JTextField} stores text, but a robust form needs raw text, semantic value, parse state, validation problems, dirty state, and commit policy.\\
Swing connection & Field bindings, validation tooltips, validation borders, dialog validation, problem summaries, and problem-focus resolution.\\
Companion example & \swingIntegrationExample.\\
\end{chaptermap}

\section{Forms are not bags of components}

A form is a contract for editing a value. A Swing panel is one possible visual
arrangement of that contract. Confusing the two is the beginning of many Swing
validation problems. A \api{JTextField} can store the characters the user typed.
It cannot, by itself, say which typed value is still semantically valid, whether
the user has touched the field, whether the field differs from its original
baseline, whether a parse failure should block commit, which resource key names
the validation message, or which component should receive focus when the user
presses OK.

\termdef{Form model}{in Corusco} is the UI-independent object that owns editable
field state, validation problems, reset behavior, dirty-state baselines, and
commit to a result. A form model can be bound to Swing, tested without Swing, or
used by generated companions.

\begin{lstlisting}[caption={The core form contract.}]
public interface FormModel<R> {
    ProblemSet problems();
    boolean isCommittable();
    void reset();
    void acceptCurrentValues();
    R toResult();
}
\end{lstlisting}

The contract is small, but it asks the right questions. Can the form currently
commit? What problems explain the answer? What is the baseline? What is the
result? Swing components are not mentioned.

\begin{figure}[htbp]
\centering
\begin{minipage}[t]{0.48\linewidth}
\includegraphics[width=\linewidth]{\textValidationFigure}
\end{minipage}\hfill
\begin{minipage}[t]{0.48\linewidth}
\includegraphics[width=\linewidth]{\dialogValidationFigure}
\end{minipage}
\caption{Form state appears visually as raw text, parse state, validation summaries, borders, and dialog commit policy. The model owns the facts before the Swing components present them.}
\label{fig:forms-validation-screens}
\end{figure}

\textbf{Parsing is not validation.}

\termdef{Parsing}{in form editing} asks whether raw text can be converted into a
semantic Java value. \termdef{Validation}{in form editing} asks whether a
semantic Java value is acceptable for the business rule. These are different
questions and should be modeled separately.

A user typing a credit limit may temporarily type \api{12,}. In some locales
that text may be incomplete. The last semantic value might still be
\api{12}. The field should display exactly what the user typed, but the form
should not corrupt the last meaningful value. A validation rule such as "credit
limit must not exceed 10000" should run after a decimal value exists, not while
the text is unparsable.

\begin{pitfallbox}{The Exception-shaped Form}
Normal user input mistakes are not exceptional program failures. Treat parse
and validation failures as data. Reserve exceptions for broken converters,
services, or programmer errors.
\end{pitfallbox}

\textbf{Raw text and semantic value.}

\api{TextFieldModel<O,T>} exists because a text field has two simultaneous
truths. The raw text is what the user sees and edits. The semantic value is the
last parsed Java value. On parse success, both move forward. On parse failure,
raw text moves forward, the semantic value remains protected, parse state
records failure, and a parse problem targets the field.

\begin{lstlisting}[caption={Text editing preserves invalid intermediate input.}]
TextFieldModel<CustomerEdit, BigDecimal> creditLimit =
        new TextFieldModel<>(
                CustomerFields.CREDIT_LIMIT,
                new BigDecimal("10.00"),
                Converters.bigDecimal(EmptyTextPolicy.REJECT));

creditLimit.setRawText("12,", ChangeOrigin.USER);

String visibleText = creditLimit.rawText().value();
BigDecimal semanticValue = creditLimit.value();
ParseState<BigDecimal> parseState = creditLimit.parseState().value();
ProblemSet parseProblems = creditLimit.problems();
\end{lstlisting}

The important observation is that none of these values is redundant. Raw text is
needed for display. Semantic value is needed for business logic. Parse state is
needed for explanation. Problems are needed for presentation, summaries, tests,
and commit policy.

\textbf{ParseState.}

\api{ParseState<T>} is the explicit state of conversion. A parsed state carries
the raw text and the parsed value. A failed state carries raw text, the previous
semantic value, and a parse problem. This avoids a common Swing anti-pattern:
storing unparsable text in a domain field merely because the text field changed.

\begin{lstlisting}[caption={Interpreting parse state.}]
return switch (field.parseState().value()) {
    case ParseState.Parsed<BigDecimal> parsed ->
            "Parsed " + parsed.value();
    case ParseState.Failed<BigDecimal> failed ->
            failed.problem().message();
};
\end{lstlisting}

Use this interpretation at the presentation boundary. The form model should own
the state; the view should decide how to display it.

\textbf{FieldModel.}

\api{FieldModel<O,T>} models a field whose current value is already semantic.
It owns a typed \api{FieldKey}, a current value, dirty state, touched state, and
the original dirty-state baseline. It is suitable for check boxes, combo boxes,
radio groups, and fields whose component already edits the semantic type.

\begin{lstlisting}[caption={A semantic field model for a checkbox.}]
FieldModel<CustomerEdit, Boolean> active =
        new FieldModel<>(CustomerFields.ACTIVE, true);

active.setValue(false, ChangeOrigin.USER);
boolean dirty = active.isDirty();
boolean touched = active.isTouched();
\end{lstlisting}

The distinction between dirty and touched is essential. Dirty means the current
semantic value differs from the baseline. Touched means the user interacted with
the field. A field can be touched but not dirty if the user changes a value and
changes it back. A field can be dirty because a model load changed it, even
though the user did not touch it.

\begin{principlebox}{Dirty Is Not Touched}
Dirty protects data from accidental loss. Touched protects the user from
premature scolding. Use both.
\end{principlebox}

\textbf{Baselines, reset, and accept.}

Every editable form needs a baseline. Reset returns fields to the baseline.
Accepting current values changes the baseline to the current values. These
operations are not the same as saving to a database, although a save operation
often calls \api{acceptCurrentValues} after it succeeds.

\begin{lstlisting}[caption={Typical save baseline flow.}]
if (form.isCommittable()) {
    CustomerEdit result = form.toResult();
    repository.save(result);
    form.acceptCurrentValues();
}
\end{lstlisting}

This sequence says: produce a committed result, save it, and only then treat the
current field values as clean. If saving fails, dirty state should remain dirty.

\section{Problems as data}

\termdef{Problem}{in Corusco} is an immutable description of parse, validation,
or feedback trouble. It has a stable code, severity, target, source, and
diagnostic message. \api{ProblemSet} is an immutable deterministic collection of
problems. Adding or filtering returns a new set.

\begin{lstlisting}[caption={A validation problem with stable identity.}]
Problem requiredName = Problem.validation(
        ProblemCode.of("customer/name/required"),
        ProblemSeverity.ERROR,
        ProblemTarget.field(CustomerFields.NAME),
        "Name is required");

ProblemSet problems = ProblemSet.of(requiredName);
\end{lstlisting}

The message is diagnostic text. User-facing localization should use resource
descriptors or application presentation code. The stable code and target are the
important contract. They let tests assert the problem without depending on a
translated sentence.

\begin{coruscobox}{Structured Failure}
Corusco forms return problem sets for ordinary user mistakes. They do not force
normal validation paths into exceptions, labels, or component client properties.
\end{coruscobox}

\textbf{Targets.}

A problem target answers "where does this problem belong?" Common targets are a
form, a field, or a table cell. Targeting matters because the same problem can
be displayed in several places: an inline border, a tooltip, a dialog summary,
a status line, or a test assertion. If the problem has no target, the UI must
guess.

For beginners, the practical rule is simple. If a validation rule concerns one
field, target that field. If it concerns several fields, target the form or the
most useful recovery field and include a summary. If it concerns a table cell,
target the row and column identity, not the current view coordinates.

\textbf{Severity.}

Severity is not decoration. It affects commit policy and presentation priority.
An error usually blocks commit. A warning may allow commit but should be visible.
Informational problems may explain formatting or policy. \api{ProblemSet} can
answer whether it contains errors and can return problems ordered by severity.

\begin{lstlisting}[caption={Committability based on problem severity.}]
public boolean isCommittable() {
    return !problems().hasErrors();
}
\end{lstlisting}

Do not encode severity only in message text. A string that begins with "Warning"
is not a model.

\textbf{Validation rules.}

Validation should be deterministic and side-effect-light. A rule reads semantic
model state and returns a \api{ProblemSet}. It should not update labels, show
dialogs, or mutate unrelated fields. Those are presentation decisions.

\begin{lstlisting}[caption={Conceptual validation rule.}]
FieldValidator<String> required =
        Validators.required(ProblemCode.of("customer/name/required"));

ProblemSet nameProblems = required.validate(CustomerFields.NAME, name.value());
\end{lstlisting}

The exact helper used by generated code may be richer, but the principle is the
same: validation returns structured problems. It does not push messages into
components.

\textbf{Rule sets and dependencies.}

A form-level rule often depends on specific fields. A date range rule depends on
start date and end date. A credit rule may depend on customer type and credit
limit. Corusco \api{RuleSet} can collect validation rules with dependency and
timing metadata so generated or handwritten form models know which rules to
refresh.

This is a maintenance advantage. In ordinary Swing code, a document listener on
one field may remember to validate a second field. In generated or model-backed
code, the dependency is declared where the rule is defined.

\begin{migrationbox}{Validation Dependencies}
When a listener in field A updates the error label for field B, write down the
rule dependency. That dependency belongs in validation metadata or a form model,
not in the document listener.
\end{migrationbox}

\textbf{Plain Swing baseline.}

A small Swing form often starts like this:

\begin{lstlisting}[caption={Component-owned parsing and validation.}]
saveButton.addActionListener(event -> {
    try {
        BigDecimal limit = new BigDecimal(limitField.getText());
        if (limit.compareTo(BigDecimal.ZERO) < 0) {
            errorLabel.setText("Limit must not be negative");
            return;
        }
        save(new CustomerEdit(nameField.getText(), limit));
    } catch (NumberFormatException ex) {
        errorLabel.setText("Limit is not a number");
    }
});
\end{lstlisting}

This code may be acceptable for a throwaway dialog. But it does not let the
screen show parse failure before Save. It does not expose dirty state. It does
not make problems testable without clicking. It does not distinguish diagnostic
messages from localized UI text. It does not preserve the last valid semantic
value.

\section{Corusco treatment}

The Corusco version moves parsing, semantic value, dirty state, touched state,
and problems into models:

\begin{lstlisting}[caption={Model-owned form state.}]
final class CustomerEditForm implements FormModel<CustomerEdit> {
    final TextFieldModel<CustomerEdit, String> name;
    final TextFieldModel<CustomerEdit, BigDecimal> creditLimit;
    final FieldModel<CustomerEdit, Boolean> active;

    @Override
    public ProblemSet problems() {
        return name.problems()
                .addAll(creditLimit.problems())
                .addAll(validateSemanticRules());
    }

    @Override
    public boolean isCommittable() {
        return !problems().hasErrors();
    }
}
\end{lstlisting}

The Swing panel binds text fields and controls to this model. The dialog OK
button observes \api{isCommittable} or a derived value. The tooltip policy reads
problem sets. Tests mutate the model directly.

\textbf{Active editor commit.}

Swing tables and some formatted controls have active editors. Pressing OK in a
dialog must first commit the active editor; otherwise the model may contain the
previous value while the user sees a newer value in the editor. This is a classic
Swing bug. It appears as "OK ignored my last edit".

Corusco dialog and binding helpers should make active-editor commit part of the
standard workflow. A form model alone cannot solve it because the edit may still
be in a Swing editor component. The boundary must be explicit: commit component
editor, update model, validate, then produce result.

\textbf{Async validation.}

Some validation requires a service: checking that a customer name is unique,
that a remote code exists, or that the current user may perform an operation.
This is \termdef{asynchronous validation}. It should not block the EDT, and it
must handle stale results.

\api{AsyncFieldValidation} binds a field value, a task service, and an async
validator. It tracks generations so a late result for an old value cannot
overwrite the current value's problems. It exposes problems as a readable value.

\begin{lstlisting}[caption={Conceptual async validation shape.}]
AsyncFieldValidation<CustomerEdit, String> uniqueName =
        AsyncFieldValidation.bind(
                CustomerFields.NAME,
                nameField.semanticValue(),
                taskService,
                (key, value, cancellation) -> repository.validateName(value));

ReadableValue<ProblemSet> asyncProblems = uniqueName.problems();
\end{lstlisting}

The exact field accessor depends on the model. The policy does not: launch work
away from the EDT, cancel or ignore stale work, and publish structured problems.

\begin{pitfallbox}{Stale Validation}
If the user types "Ada" and then "Adal", the result for "Ada" must not overwrite
the current field state when it completes late.
\end{pitfallbox}

\textbf{Dialogs and commit policy.}

The form chapter is the core counterpart to the dialog chapter. A dialog should
not ask each component whether OK may proceed. It should ask the form model
whether a result can be produced. If the answer is no, it should display or focus
the most useful problem. If the answer is yes, it should call \api{toResult}.

A good OK flow is:

\begin{enumerate}
\item Commit the active Swing editor.
\item Refresh or observe current validation state.
\item If blocking problems exist, display and focus them.
\item If no blocking problems exist, call \api{toResult}.
\item Close or continue according to dialog policy.
\end{enumerate}

This flow is boring. That is its virtue. Every dialog should not invent it
again.

\textbf{Testing forms.}

Test form models before testing Swing bindings. A model test should mutate raw
text, semantic fields, and validators, then assert raw text, parse state,
problems, dirty state, touched state, committability, and committed result.

\begin{lstlisting}[caption={Testing invalid intermediate input.}]
form.creditLimit.setRawText("bad", ChangeOrigin.USER);

assertThat(form.creditLimit.rawText().value()).isEqualTo("bad");
assertThat(form.creditLimit.value()).isEqualByComparingTo("10.00");
assertThat(form.creditLimit.parseState().value())
        .isInstanceOf(ParseState.Failed.class);
assertThat(form.problems().hasErrors()).isTrue();
\end{lstlisting}

The Swing test can then be narrow: typing in the field updates the model; problem
state decorates the component; pressing OK follows dialog policy.

\textbf{Worked example.}

\begin{examplebox}{Forms}
\swingIntegrationExample{} includes form and binding code that can be read beside
this chapter. The core lesson is not the visual layout; it is the separation
between raw text, semantic value, problems, and dialog behavior.
\end{examplebox}

\textbf{Review notes.}

Do not let validation disappear into labels. Do not let parsing throw ordinary
user mistakes into exception control flow. Do not let dirty state mean whatever
the current listener happens to update. A form model is the place where these
concepts become named and testable.

The first design question for a form is what result it promises to produce. A
dialog that edits a customer record, a preferences panel that writes settings,
and a search form that produces criteria are all forms, but they do not promise
the same kind of result. Naming the result keeps the field models honest. Each
field exists because it contributes to that result, explains why the result
cannot be produced yet, or records user intent that must be carried into the
workflow.

This is why raw text and semantic value must remain separate in text fields. The
user is allowed to type an intermediate string that is not yet a value. A
decimal field may temporarily contain a minus sign, a date may be halfway typed,
and a required field may be empty while the user is moving through the dialog.
If the model stores only the semantic value, these intermediate states are lost.
If the model stores only text, command logic and domain code must parse again and
again.

Validation should report structured facts rather than perform immediate UI
theatre. A red border, tooltip, summary row, disabled OK button, and status
message are presentations of problems; they are not the problems themselves.
Once problems have severity, code, target, and text resource, the UI can choose
the right presentation for the current context. The same problem may appear as a
field decoration in a dialog and as a row-level message in a review table.

Severity is an agreement with workflow. An error blocks result creation. A
warning may allow OK but should remain visible. Information may guide the user
without affecting committability. If severity is hidden inside label colors,
tests cannot assert policy and dialogs cannot behave consistently. A
\api{ProblemSet} makes the policy explicit: which problems block, which merely
decorate, and which target deserves focus first.

Dirty state is often confused with invalidity. A pristine form may be invalid
because required data is absent in the original object. A dirty form may be
valid because the user fixed the issue. A touched field may not be dirty because
the user restored the original value. Keeping these states separate prevents
annoying cancel prompts, premature validation displays, and incorrect Save
enablement.

Touched state exists for kindness. It lets a form delay some messages until the
user has interacted with a field while still keeping the full validation model
available to OK and tests. A strict accounting screen may show all errors from
the beginning. A friendly data-entry dialog may initially show only blocking
summary state, then decorate individual fields after they are touched. The model
can support both policies because touched state is not the same as validation.

Cross-field validation should be written as a rule over the form, not as a
listener attached to whichever component happened to be edited first. A credit
limit may depend on customer status. A date range depends on both dates. A
shipping option depends on address and account permissions. The rule belongs to
the form model because the dependency is semantic; component listeners are only
one way the relevant values change.

Asynchronous validation adds another boundary. The form should expose that a
check is running and should remember which value the check belongs to. The
repository should perform the slow work. The task service should keep work off
the EDT. The callback should publish problems only if the value is still
current. Mixing those responsibilities into a document listener produces exactly
the stale-result bug that users experience as validation flicker.

Dialog OK is a commitment protocol, not just a button press. The UI must commit
active editors, update models, ask the form whether the result is available,
focus or summarize blocking problems, and only then create the result. This
sequence should be boring enough to reuse. A form model supplies the stable
facts; the dialog shell supplies the Swing ceremony.

The best model tests read like small editing stories. Start from an original
record. Type bad text. Observe that the raw text changed while the semantic
value did not. Type good text. Observe that the parse problem disappeared and
the value changed. Reset the form. Observe dirty and touched state. Accept the
baseline. Observe that the current value is no longer dirty. Those stories catch
more defects than visual tests that merely click OK once.

The Swing binding tests can then stay narrow. They should prove that a
\api{JTextField} edits the correct \api{TextFieldModel}, that problem state
reaches the expected decoration, and that dialog buttons observe
committability. They need not retest every validation rule through mouse and
keyboard input because the rules already live in the model. This division keeps
the test suite fast and the screen behavior credible.

The final form-design check is whether every visible validation behavior has a
model fact behind it. A red border, disabled OK button, warning icon, tooltip,
summary row, and dirty marker should all be explainable from field models,
problem sets, touched state, dirty state, or busy state. If a decoration cannot
be traced back to a model fact, it will be hard to test and harder to keep
consistent across dialogs.

\section{Designing the result contract}

A form should be designed from the result backward. This sounds obvious, yet
many Swing forms are designed from components forward: place a label, place a
text field, add another field, add OK, and later decide what object should be
created. That order encourages the panel to become the model. A result-first
design asks what the user is trying to commit, which editable facts contribute
to that result, which facts merely guide editing, and which facts are visual
support. The answer may be an immutable record, a command request, a search
criteria object, a preferences delta, or no single object at all. The form
model's type parameter is the first sentence of that answer.

The \api{FormModel<R>} contract is deliberately small because the result
contract should not be hidden under ceremony. \api{problems} explains the
current obstacles and feedback. \api{isCommittable} answers whether a result is
available. \api{reset} restores the current baseline. \api{acceptCurrentValues}
changes the dirty baseline without changing the represented result.
\api{toResult} materializes the committed value and must fail when the form is
not committable. A small contract is not a small idea. It is a discipline: the
form is either able to produce its promised result, or it can say why not.

\api{AbstractFormModel} gives this discipline a useful default shape. A
handwritten or generated subclass registers its \api{FieldModel} and
\api{TextFieldModel} instances, implements \api{createResult}, and optionally
overrides \api{validationProblems}. The base class aggregates parse problems,
adds semantic validation problems, resets fields, accepts current values as the
new baseline, and blocks \api{toResult} while an error remains. The base owns
registration lists, not Swing components. This is important: the form model can
be tested as a model and then bound to one or several Swing views.

\begin{lstlisting}[caption={A result-first form model skeleton.}]
final class CustomerEditFormModel
        extends AbstractFormModel<CustomerEdit> {

    final TextFieldModel<CustomerEdit, String> name =
            register(new TextFieldModel<>(
                    CustomerFields.NAME, original.name(), Converters.string()));

    final TextFieldModel<CustomerEdit, BigDecimal> creditLimit =
            register(new TextFieldModel<>(
                    CustomerFields.CREDIT_LIMIT,
                    original.creditLimit(),
                    Converters.bigDecimal(EmptyTextPolicy.REJECT)));

    final FieldModel<CustomerEdit, Boolean> active =
            register(new FieldModel<>(
                    CustomerFields.ACTIVE, original.active()));

    @Override
    protected CustomerEdit createResult() {
        return new CustomerEdit(
                name.value(),
                creditLimit.value(),
                active.value().value());
    }
}
\end{lstlisting}

The skeleton is intentionally ordinary. The result is created from semantic
values, not from Swing components. The text fields preserve raw editing state,
but \api{createResult} uses the semantic values only after committability has
been checked. A generated form model follows the same rule: generated fields,
descriptors, validators, and result creation exist so the result contract is
stable and inspectable.

The baseline is part of the result contract because editing is temporal. A
customer editor opens with one customer value. The user changes fields. The
user may reset, cancel, save, or keep editing after a save failure. Without a
baseline, the form cannot answer "what would reset do?" or "is it safe to close
without confirmation?" A baseline is not necessarily the database value. In a
wizard, the baseline may be the state at the beginning of the current page. In
a preferences dialog, it may be the values when the dialog opened. In an
import-review screen, it may be the parsed file state after automatic cleanup.
The form model should make that choice explicit.

\api{reset} and \api{acceptCurrentValues} are therefore workflow operations.
\api{reset} says that current edits are abandoned in favor of the baseline.
\api{acceptCurrentValues} says that current semantic values are now considered
the baseline. A successful save commonly calls \api{acceptCurrentValues}; a
failed save should usually not. A successful Apply operation may accept the
baseline while leaving the dialog open. A successful OK operation may not need
to accept because the dialog closes, but accepting can still be useful if the
same model survives in a larger screen. These choices should be written in the
dialog or presenter, not smuggled into a button listener.

Dirty state is only meaningful relative to that baseline. A field is dirty when
its current semantic value differs from the original value remembered by its
field model. A text field whose raw text is unparsable can be touched and can
have parse problems even while its semantic value remains unchanged. This can
surprise code that treats dirty as "the characters in the text box changed."
Corusco's split is more useful: dirty says the semantic value differs from the
baseline; touched says the user interacted with the field; parse state says
whether raw text currently represents a semantic value.

Consider a credit-limit field whose baseline is \api{100.00}. The user selects
the field, deletes the text, types \api{-}, then types \api{-5}, then changes
it back to \api{100.00}. During the first edit the field is touched. During the
invalid intermediate text, the raw text changed, the semantic value remains
protected, and a parse problem may exist. During \api{-5}, parsing may succeed
and semantic validation may reject the value. At the end, the field may still
be touched but no longer dirty. A single boolean named \api{changed} cannot
describe this story. The form model can.

Touched state is not merely a convenience for prettier validation. It is a
policy hook. A form may initially contain a required but empty field. Should
the field show a red border immediately? In an accounting screen, perhaps yes:
the user opened the form to fix known errors. In a small creation dialog,
perhaps no: the user has not yet had a chance to type. Touched state lets the
presentation layer distinguish "this fact blocks commit" from "this fact should
be loudly decorated right now." The problem exists in the model either way.
The visual policy can be kind without lying.

The result contract also prevents accidental partial commits. A form whose
\api{isCommittable} returns false should not return a half-formed result from
\api{toResult}. This is not only defensive programming. It protects maintainers
from ambiguous workflow. If a caller wants a partial object for preview, name a
different method or a different model. The committed result should mean that
the form's blocking problems were absent at the time the result was created.
The exception raised by \api{toResult} on an uncommittable form is the model
refusing to pretend.

\begin{detailbox}{Result Creation Is Not Validation}
The \api{createResult} method of an \api{AbstractFormModel} should normally
assemble already-validated semantic values. It should not discover ordinary
parse or validation failures for the first time. Those failures belong in
\api{problems} so the UI can explain them before commit.
\end{detailbox}

The result-first habit also helps when there is no classic OK button. A search
panel may produce criteria on every change. A filter form may produce a
predicate. A preferences page may produce a diff. A bulk-edit panel may produce
a command request that applies to selected rows. These are still forms in the
chapter's sense if they own editable state, problems, baseline, and result
policy. They may not use \api{JDialog}, but they still benefit from a
\api{FormModel}-shaped contract.

\section{Conversion as a boundary, not a nuisance}

Text is deceptively simple. A \api{JTextField} contains a \api{Document}; a
\api{Document} stores characters and participates in input methods, undo,
filters, caret movement, and document events. That does not mean the
application has a value. The application has a value only after conversion has
interpreted the characters according to the field's type, locale, empty-text
policy, and formatting rules. Conversion is the boundary where editing becomes
meaning.

A converter has two directions. Formatting turns a semantic value into text.
Parsing turns text into a semantic value or a parse failure. Both directions
should be considered when designing a field. A date may format as
\api{2026-06-14} in one screen and as a localized date in another. A decimal
may accept a comma in a locale-aware converter. An optional field may format a
\api{null} value as an empty string. A required field may still allow an empty
string as raw text while producing a parse or validation problem. These
policies should be explicit because users edit text, while the application
commits values.

\api{TextFieldModel} exists to keep this boundary honest. It owns raw text,
parse state, parse problems, dirty state, touched state, and a semantic
\api{FieldModel}. When \api{setRawText} succeeds, the semantic field advances
and parse problems are cleared. When parsing fails, raw text still advances,
the semantic value remains the previous good value, touched state is marked,
and a parse problem targets the field. This gives the user freedom to edit
without forcing the rest of the application to defend itself against corrupted
semantic state.

The previous semantic value in \api{ParseState.Failed} is more than a nice
detail. It makes possible a stable preview, a stable dependent field, and a
stable command rule while the user repairs text. Suppose a tax-rate field
drives a calculated preview. If the user temporarily types \api{.}, the
preview can either hold the last meaningful calculation or show that the input
is not currently usable. What it should not do is treat \api{.} as zero,
overwrite the model with nonsense, and then run business rules against the
nonsense. The failed parse state gives the presenter a truthful choice.

Parsing and validation should not duplicate each other. If text cannot be
parsed as a date, a date-range validator should not also report that the date
is outside the range. The user needs one precise first problem: the text is not
a date. \api{RuleSet.Builder.field} encodes this policy. A rule for a parsed
text field checks the field's parse problems and skips semantic validation when
parse errors are already present. This is not hiding errors. It is preserving
the order of explanation. First make the text a value; then ask whether the
value is acceptable.

\begin{lstlisting}[caption={RuleSet skips semantic validation after parse failure.}]
RuleSet<CustomerEditFormModel> rules =
        RuleSet.<CustomerEditFormModel>builder()
                .field(CustomerFields.CREDIT_LIMIT,
                        form -> form.creditLimit,
                        Validators.decimalRange(
                                BigDecimal.ZERO,
                                new BigDecimal("10000.00"),
                                "Credit limit is outside the allowed range"))
                .build();
\end{lstlisting}

In this example the range rule is meaningful only when a decimal exists. If the
raw text is \api{abc}, the \api{TextFieldModel} owns the parse problem and the
range rule stays silent. If the raw text parses to \api{-1}, the parse state is
fine and the range rule may report a validation problem. The distinction makes
error summaries shorter and user repair paths clearer.

Empty text deserves special treatment because it is both common and ambiguous.
An empty string in a required field may mean "the user has not entered the
value yet." An empty string in an optional numeric field may mean \api{null}.
An empty string in a search field may mean "no filter." An empty string in a
code field may be syntactically invalid. Do not make one global converter
policy do all jobs. The converter and validator together should explain
whether empty raw text is parseable, whether the resulting semantic value may
be absent, and whether absence blocks commit.

Formatted values are another source of accidental hostility. A converter that
reformats raw text on every keystroke may fight the user's cursor. A currency
field that inserts grouping separators while the user types can move the caret
and make editing unpleasant. A form model should preserve raw text; a binding
or editor policy can decide when to reformat. Common policies are: format on
reset, format on commit, format when focus leaves the field, or format only
when model changes originate from code rather than user typing. The important
point is not which policy wins universally. The important point is that the
policy is explicit.

Input methods make this caution stronger. International text input may involve
composition before committed characters appear in the document. Key listeners
that try to parse individual keystrokes are brittle. Document and binding
updates should treat the field as text edited by Swing's text machinery.
Parsing raw text after document changes is reasonable; assuming every key press
is a final character is not. The accessibility and input chapter returns to
this point, but form code should already avoid shortcutting text components.

The converter itself should be deterministic and side-effect-light. It should
not query a remote service, update labels, or mutate other fields. If parsing a
value requires a remote lookup, then the local parse result is probably an id,
code, or raw token, and the remote lookup is validation or enrichment. Keeping
conversion local keeps editing responsive and testable. A slow converter is
just another way to block the EDT.

\begin{pitfallbox}{Formatting as Validation}
Do not silently repair user text and call that validation. Trimming surrounding
space may be harmless in one field and destructive in another. Normalization is
a domain policy. If it changes meaning, make it visible in the model or result.
\end{pitfallbox}

There are valid cases for strict input filters. A field that accepts only a
single digit, a mask that represents a fixed code, or a custom editor with a
small finite vocabulary may reject impossible edits before they reach the
model. But filters should be used carefully. A filter that rejects every
intermediate invalid state can make ordinary editing impossible. The user must
be able to delete the old value before typing the new one. The
\api{TextFieldModel} approach is usually kinder: accept raw text, preserve the
last semantic value, and report parse state.

\section{Problems, targets, and presentation policy}

A problem is a routed fact. It is not merely a message. Corusco's
\api{Problem} record contains a stable code, severity, target, source, and
diagnostic message. The code lets tests, generated code, and resource lookup
identify the problem without depending on wording. The severity lets workflow
decide whether commit is blocked. The target tells the UI where recovery
belongs. The source distinguishes parse, validation, and other producers. The
message is useful for tests and logs, but user-facing localization belongs in
resources and presentation policy.

\api{ProblemSet} is immutable and deterministic. Adding problems returns a new
set. Filtering returns a new set. Iteration preserves insertion order.
Severity-ordered views are stable for problems of equal severity. The class
does not deduplicate. That last rule is deliberate. If two validators report
the same problem, the transport should not silently decide whether that is a
bug, a merge opportunity, or two independent findings. Presentation code can
merge or suppress duplicates when it has the proper context.

Targeting is the difference between feedback and noise. A problem targeted at
\api{ProblemTarget.field(CustomerFields.NAME)} can decorate the name field,
appear in a summary link that focuses the name field, contribute to the name
field's accessible description, and be asserted by a model test. A form-wide
problem can appear above the form or in a status area. A table-cell problem can
decorate a row and column identity without relying on current view coordinates.
A component target can be used when the relevant UI element is not naturally a
field. Without targets, every presentation is guesswork.

Problem targets should use stable identities, not visible text. A field label
may be localized. A table row may move. A table view column may be reordered.
A component may be nested in another panel. A \api{FieldKey},
\api{ComponentKey}, row identity, or column identity survives those changes.
This is the same identity discipline used elsewhere in the book: what the user
sees is presentation; what the program routes through should be stable.

Severity is equally practical. \api{ProblemSeverity.ERROR} normally blocks
commit. \api{ProblemSeverity.WARNING} and \api{ProblemSeverity.INFO} normally
do not, unless a caller applies stricter policy. The declaration order is the
natural order from least severe to most severe, which means severity filters
and sorting have a predictable foundation. The UI should not infer blocking
policy from colors, icons, or message prefixes. A yellow icon is presentation;
a warning severity is model data.

A form may have several presentation channels for the same problem. Inline
decoration helps the user see where to fix a field. A tooltip can provide a
short explanation. A summary helps in long forms and supports keyboard
navigation. Accessible descriptions make problems available to assistive
technology. Status text can explain why a command is disabled. The problem
model should feed these channels; none of them should be the original source
of truth. Otherwise they will drift.

The display policy should answer three questions. Which problems are shown
immediately? Which problems block commit? Which problem receives focus after a
failed commit attempt? The answers need not be identical. A strict form may
show all errors immediately. A gentle form may show field-level decorations
only after touch, while still blocking OK on untouched required fields. A
failed OK may focus the most severe field problem or the first problem in
reading order. A warning may remain in the summary after a successful parse but
not block OK. These policies belong in the dialog or binding layer over the
\api{ProblemSet}.

\begin{lstlisting}[caption={Selecting a focus target from structured problems.}]
Optional<Problem> focusProblem = form.problems()
        .bySeverityDescending()
        .stream()
        .filter(problem -> problem.target()
                instanceof ProblemTarget.Field<?, ?>)
        .findFirst();
\end{lstlisting}

The snippet is not a complete focus resolver, because real focus policy should
also consider field visibility, enabled state, current page, table row
visibility, and possibly touched state. It shows the key point: once problems
have targets and severities, focus policy can be written as policy. Without
targets, focus policy becomes string matching or a list of special cases.

Messages need a similar policy. The \api{Problem} message is diagnostic text.
It is useful for tests, logs, and early examples. A production application
usually wants resource-backed, localized text. The stable problem code and
target can select a resource. The field descriptor can supply a label resource.
The presentation layer can combine them: "Customer name is required" may come
from a validation message template and a field label. Keeping this composition
out of the validator makes validators reusable and testable.

The problem source is a small but useful diagnostic. A parse problem and a
validation problem can have the same target but different remedies. A parse
problem asks the user to make text into a value. A validation problem asks the
user to make the value acceptable. A task or service problem may ask the user
to retry, wait, or change connectivity. Source also helps presentation avoid
duplicating messages. If parse failed, semantic validation for that field
should usually be skipped. If async validation failed unexpectedly, the UI may
show a different tone from an ordinary "name already exists" validation result.

\begin{coruscobox}{Problem Sets Are Presentation-neutral}
Corusco does not say whether a problem appears as a red border, icon, tooltip,
summary row, accessible description, or disabled reason. It gives those
presentations one structured fact to observe.
\end{coruscobox}

Problem ordering should not be accidental. \api{ProblemSet} preserves insertion
order, so rule registration order can express reading order or workflow
priority. The severity view then places severe problems first while retaining
registration order within the same severity. This lets a form produce stable
summaries. Stable summaries matter for screenshots, tests, and user trust. A
summary that changes order on every validation pass makes the form feel
unstable even if the underlying problems are correct.

The route from problem to component should be reversible enough for tests.
Given a field key, a binding should know which component displays the field.
Given a problem target, a focus resolver should know which component can
receive focus. Given a component, diagnostics should be able to report the
field or component key. Corusco generated view contracts and component keys
help here, but the principle is independent: validation feedback should be
routed by identity, not by searching visible labels.

Some problems are intentionally not field-local. A date range may target the
form because neither date alone is wrong. A summary may also choose a recovery
field, usually the field the user should edit first. A table import may target
a row when the entire row conflicts with existing data. A connection failure
may target the form or command. Do not force every problem into a field. Do
ensure every problem has a target whose presentation policy is known.

\section{Validation lifecycle and asynchronous checks}

Validation has a lifecycle. It is not a single call hidden behind OK. A form may
validate on every change, after focus loss, before commit, after a background
service returns, after a dependent field changes, or after a table row is
edited. Different rules belong at different moments. Corusco's
\api{ValidationTiming} names two common moments: \api{ON\_CHANGE} and
\api{ON\_COMMIT}. The core does not schedule those moments; it lets form models,
generated plans, and Swing bindings choose which rules to run.

Immediate validation is useful when feedback helps the user continue. A
required field, length rule, range rule, or format-sensitive dependency may be
cheap enough to run whenever relevant values change. Commit validation is
useful when the rule is expensive, cross-cutting, or too noisy during ordinary
typing. A rule that checks whether a date range overlaps an existing booking
may not need to run on every keystroke. A rule that enforces final agreement
with several fields may be better saved for OK or Apply.

\api{RuleSet} gives validation rules a typed dependency list. A rule that
depends on start date and end date should say so. A rule that depends on
customer type and credit limit should say so. This avoids the old Swing pattern
where a listener on start date updates an error label for end date because the
developer remembered the relationship in one place but not another. Dependency
metadata is executable documentation: it tells generated or handwritten code
which rules are affected by a changed field.

\begin{lstlisting}[caption={A cross-field rule with explicit dependencies.}]
RuleSet<CustomerEditFormModel> rules =
        RuleSet.<CustomerEditFormModel>builder()
                .form(List.of(CustomerFields.START_DATE,
                              CustomerFields.END_DATE),
                        form -> validateDateRange(
                                form.startDate.value(),
                                form.endDate.value()))
                .build();
\end{lstlisting}

The validator in the example reads semantic values. It should not inspect text
fields, labels, or borders. If either text field has parse problems, the form
model can decide whether the cross-field rule should skip, return no
additional problems, or report a form-level problem about incomplete input. The
important point is that this is a semantic rule over the model, not a side
effect of editing one component.

Validation should be side-effect-light. A field validator returns a
\api{ProblemSet}. A form validator returns a \api{ProblemSet}. They should not
show message boxes, mutate unrelated fields, enable buttons, or install
listeners. Those actions belong to presenters and bindings. Side-effect-light
validators are easy to test. They are also easier to generate because the
annotation processor can wire rules without guessing which UI side effects a
lambda performs.

The built-in \api{Validators} helpers follow this shape. \api{required}
returns a problem for \api{null} and blank strings. Range validators accept
missing optional values and report one field-targeted problem when a present
value is outside bounds. Regex validators treat absent strings as acceptable
unless combined with \api{required}. Date validators use a \api{Clock}, which
keeps tests deterministic. These details matter because validators are
reusable model code, not ad hoc UI fragments.

Clock injection deserves a small pause. A rule such as "date must be in the
future" can become nondeterministic if it reads the system clock directly. A
test written at midnight or across a time-zone boundary can fail for reasons
unrelated to the form. A validator that accepts \api{Clock} says time is an
input. The book uses deterministic examples for screenshots; validators should
show the same discipline in model tests.

Asynchronous validation is a different lifecycle. A rule that queries a
repository, server, cache, or filesystem should not block the EDT. It should
also not publish stale results. The user can type faster than the service can
answer. If the form validates \api{Ada}, then the user types \api{Adal}, then
the service returns the old \api{Ada} result, the old result must not overwrite
the current field's problems. This is the bug generation counters are designed
to prevent.

\api{AsyncFieldValidation} owns one asynchronous validation relationship. It
observes a semantic \api{ReadableValue}, submits work to a \api{TaskService},
exposes observable problem and busy values, and uses a
\api{GenerationCounter} so only the current generation may publish problems.
Closing the controller invalidates generations, removes the value
subscription, cancels tracked tasks, clears handles, and publishes not-busy
state. This is not merely validation code. It is lifecycle code.

\begin{lstlisting}[caption={Async validation is a lifecycle-owned relationship.}]
AsyncFieldValidation<CustomerEdit, String> uniqueName =
        AsyncFieldValidation.bind(
                CustomerFields.NAME,
                form.nameValue(),
                taskService,
                (key, value, cancellation) ->
                        repository.validateUniqueName(value, cancellation));

scope.add(uniqueName);
\end{lstlisting}

The exact accessor may differ in generated code, but the ownership rule is
stable. The async validator does slow work through the task service. The
controller owns subscription and task handles. The scope closes it with the
screen. The form or presenter observes its problems and busy state. The text
field does not own the remote query.

Busy state from async validation has presentation consequences. A field may
show a small pending indicator. A summary may say that validation is still
running. OK may be disabled while blocking async validation is current, or OK
may run a final commit validation that waits for or rechecks the service. Both
policies can be valid. What is not valid is letting a late background callback
enable OK after the user changed the field. Current generation and current
value must agree.

Async failure also needs policy. \api{AsyncFieldValidation} converts validator
exceptions and task-submission failures into a field-targeted error problem
with \api{ASYNC\_VALIDATION\_FAILED}. That default is useful because unexpected
validation failure should be visible and should usually block commit. A product
may choose a more specific message or recovery path, but the problem still has
a code, severity, target, and source. It is not a stack trace dropped into a
label.

Cancellation is cooperative. A task service can cancel a task handle, and the
validator receives a cancellation token. A validator that performs repeated
work should check the token. A validator that calls a blocking external API
may not be able to interrupt immediately, but generation checks still prevent
stale publication. Cancellation improves resource usage; generation checking
protects correctness. A robust async form wants both.

Validation timing and async validation meet at commit. Pressing OK should not
blindly trust that every async result is current unless the form's policy says
so. Some screens require all pending async checks to finish. Some screens allow
commit with warnings and perform server-side validation during save. Some
screens run a final synchronous validation on the save endpoint and translate
server responses into problems. The form chapter does not prescribe one
policy, but it insists that the policy be explicit. "Whatever callback happened
last" is not a policy.

\section{Dialogs, generated forms, and Swing bindings}

The form model is Swing-free, but a Swing form still has component mechanics.
Text fields have documents. Formatted fields may have edit/commit behavior.
Tables have active cell editors. Combo boxes have editor components. Focus can
remain in a field whose document has not yet been observed by the model. A
dialog OK handler that ignores active editors can validate stale model state
while the user sees a different value on screen. The bug is ordinary and
infuriating: "I typed a value and OK ignored it."

The correct boundary sequence is component commit, model validation, result
creation, dialog close. Component commit means asking Swing editors to push
their current editing state into the model-facing component or model adapter.
For a table, this may mean stopping cell editing. For a formatted text field,
it may mean committing or reverting the edit according to policy. For ordinary
text fields bound through documents, the model may already be current, but the
dialog should still have a standard pre-OK phase rather than a pile of special
cases.

\begin{lstlisting}[caption={A dialog OK sequence as a reusable protocol.}]
void ok() {
    activeEditorCommit.commitOrReport();
    ProblemSet problems = form.problems();
    if (problems.hasErrors()) {
        problemPresenter.showAndFocus(problems);
        return;
    }
    CustomerEdit result = form.toResult();
    closeWithResult(result);
}
\end{lstlisting}

The snippet names collaborators rather than hiding them. \api{activeEditorCommit}
knows Swing editor mechanics. The form knows problems and result creation. The
problem presenter knows visual and focus policy. The dialog shell knows how to
close. If a future bug appears, the code already tells the maintainer where to
look.

Generated Corusco forms should reinforce this division. A generated
\api{@CoruscoForm} model can register text fields and semantic fields, expose
descriptors, wire validation rules, and create the immutable record result. It
should not force a visual layout. The application panel can still be a
handwritten MigLayout panel. Generated bindings can attach model fields to the
components exposed by a view contract. This lets generation remove repetitive
wiring without turning the generator into a GUI builder.

The generated view contract is a practical compromise. Generated binding code
needs to know that there is a name text field, an active check box, a customer
type combo box, and perhaps a validation summary. The application view decides
where those components live. The same form model might appear in a compact
dialog, a wide detail pane, and a wizard page. If the generator owned layout,
those three presentations would either become uniform and awkward or require a
second set of layout escape hatches. Keeping layout handwritten preserves the
MigLayout discipline established earlier.

\begin{migrationbox}{Generated Form Adoption}
Do not introduce generation by asking it to redesign a screen. First extract
stable field keys, result type, converters, validation rules, and resource
keys. Then let generation compile those stable facts into model and binding
companions while the existing MigLayout view remains visible and reviewable.
\end{migrationbox}

Bindings should be scoped. A binding from \api{JTextField} to
\api{TextFieldModel} installs document listeners and value subscriptions. A
validation decoration binding observes problem state and changes border,
tooltip, accessible description, or icon state. A command binding observes
committability. These relationships are not free-floating. They live as long
as the screen or dialog lives. A \api{BindingScope} or lifecycle scope should
own them and close them when the view closes.

Bindings should also avoid feedback loops. A user edit changes the document.
The binding writes raw text to the field model with a user origin. The model
notifies observers. The binding may need to update the document when model
origin is reset or load, but it should not rewrite the document unnecessarily
for the same user edit. Equality checks, update guards, and change origins are
boring but essential. Without them, a form can flicker, move the caret, or
loop until the EDT is saturated.

Problem presentation bindings need restoration. If a binding changes a field's
border to show an error, it should remember the previous border and restore it
when the problem disappears or the binding closes. If a binding appends
validation text to a tooltip, it should not destroy the field's help tooltip
forever. If a binding sets accessible description, it should compose with help
and disabled reasons according to policy. Validation decoration is component
state; component state needs lifecycle.

Dirty cancel confirmation is another dialog workflow over the form model. When
the user presses Cancel or closes the window, the dialog asks whether the form
has dirty semantic values or other unsaved work. If not, it closes. If yes, it
may ask for confirmation. A touched but not dirty field should not usually
trigger "discard changes." An invalid raw text state may deserve confirmation
because the user typed something, even if no semantic value changed. This is a
policy choice, but the model gives the facts needed to make it.

Apply buttons require even more precision. OK is terminal. Cancel is terminal.
Apply is not terminal. After Apply succeeds, should the baseline be accepted?
Usually yes, because the current values have been applied. If Apply fails,
dirty state should remain. If Apply partly succeeds, the dialog needs a policy:
accept the successful fields, keep all dirty, or reload from the authoritative
state. The model cannot invent this policy. It can make the consequences
visible through baselines and problems.

Search forms are worth mentioning because they often do not look like dialogs.
A search field, date range, status combo, and Search button form an editable
contract. The result may be a search criteria object or an immediate update to
an observable list. Parse problems still matter. Dirty state may mean "criteria
changed since last search." Apply may mean "run search." Reset may mean
"restore default criteria." Treating search controls as a form avoids a tangle
of document listeners and stale result flags.

Table-cell editors are forms in miniature. A cell editor may display raw text
while the table model stores semantic values. A parse failure should not
silently corrupt the row object. A validation problem should target the cell by
row identity and column identity, not by current view row and view column.
Corusco table bindings and form models can cooperate here: the table chapter
handles row and column identity; the form chapter supplies conversion and
problem semantics.

\section{Migration and review scenarios}

Many existing Swing forms cannot be rewritten in one pass. The goal is not to
replace every listener with a framework-shaped object overnight. The goal is to
move facts to their correct owners. A good migration begins by identifying the
form result, the fields that contribute to it, the validation facts currently
hidden in UI code, and the lifecycle of installed listeners. Then extract the
most painful concept first.

If users report "Save is enabled at the wrong time," extract committability.
Create a form model or presentation value that answers whether Save can
proceed and why. Bind the button to that fact. Do not fix the bug by adding one
more listener that toggles the button after a special case. If users report
"error messages disagree," extract problems. If users report "Cancel warns
when I changed nothing," extract dirty and touched state. If users report "OK
ignored my table edit," extract active-editor commit into the dialog workflow.

\begin{lstlisting}[caption={A migration seam around legacy component state.}]
final class LegacyCustomerFormAdapter
        implements FormModel<CustomerEdit> {

    private final LegacyCustomerPanel panel;
    private final CustomerEditFormModel model;

    void pullFromComponents() {
        model.name.setRawText(panel.nameField().getText(), ChangeOrigin.USER);
        model.creditLimit.setRawText(
                panel.creditLimitField().getText(), ChangeOrigin.USER);
        model.active.setValue(
                panel.activeBox().isSelected(), ChangeOrigin.USER);
    }

    @Override
    public CustomerEdit toResult() {
        pullFromComponents();
        return model.toResult();
    }
}
\end{lstlisting}

This adapter is not the final architecture. It is a migration seam. It lets
tests begin asserting model behavior while the old panel still exists. Over
time, bindings can replace \api{pullFromComponents}, and component scraping can
disappear. A migration seam is acceptable when it moves authority toward the
model. It is not acceptable when it becomes a permanent way to hide the model
behind the panel.

The first review scenario is the integer field. The old form rejects non-digit
keystrokes. Users cannot easily type a negative number, replace all text, or
paste a value with surrounding spaces. The reviewer should ask whether the
field really needs strict character filtering or whether it should preserve raw
text and report parse state. If the domain accepts negative values, a filter
that rejects \api{-} is wrong. If the domain rejects negative values, the parse
may still accept \api{-1} and validation should report the semantic rule. This
is how a small field teaches the parsing-validation split.

The second scenario is the date range. Start date and end date each parse
independently. The range rule depends on both semantic dates. If either date
fails to parse, the range rule should not produce a misleading "start must be
before end" problem. If both dates parse and the order is wrong, the rule may
target the form or the end date according to recovery policy. The problem
summary should not depend on which document listener happened to run last.

The third scenario is an async uniqueness check. The old form starts a server
request from a document listener and writes the result into a label. Users see
the label flicker, and sometimes the label reports that a previous name is
available after they have typed a new one. The migration extracts a semantic
name value, an \api{AsyncFieldValidation}, a busy indicator, and a generation
check. The server result becomes a problem set for the current generation only.
The label becomes one presentation, not the authority.

The fourth scenario is Apply. A preferences dialog has OK, Cancel, and Apply.
The old Apply writes current component values to settings and leaves dirty
state unchanged. The user presses Apply successfully, then Cancel, and the
dialog asks whether to discard changes that were already applied. The fix is
not another boolean in the button listener. The fix is a baseline policy:
after successful Apply, accept current values as the new baseline. After failed
Apply, keep dirty state. If some settings are applied by a background task,
the baseline should change only when the task succeeds.

The fifth scenario is import review. A CSV import screen shows parsed rows in a
table and a detail form for the selected row. Parse errors in the file target
rows or cells. Edits in the detail form target fields. A duplicate check may
be async. The final Import command should ask whether blocking problems remain
in the row set and detail form. This is a larger workflow, but it is made of
the same pieces: raw text, semantic values, problem targets, severities,
baselines, selection, and task lifecycle.

The sixth scenario is generated forms. A team annotates a record and receives a
generated model, descriptors, and binding companion. The generated code looks
strange at first because it names facts that the old form kept implicit:
field keys, converter policy, validators, result creation, and view contract.
The reviewer should ask whether those generated names are stable and useful.
If they are, the generator is compiling real design. If they are awkward, the
annotations may be exposing unclear domain language that was already present
in the old form.

Review should also include screenshots, but screenshots are evidence of
presentation, not proof of model correctness. A screenshot can show that a
problem summary is readable, a tooltip is not clipped, a dialog does not jump
when validation appears, and a disabled button uses the right styling. It
cannot prove that stale async validation is impossible or that \api{toResult}
throws when errors exist. Pair screenshot examples with model tests.

The model tests should be written as editing stories. A good story names the
baseline, user edit, parse result, validation result, dirty state, touched
state, and commit result. For async validation, the story should include stale
completion. For dialogs, the story should include active-editor commit and OK
failure focus. For generated forms, the story should include annotation
processor output compiling and the generated model producing the expected
record. These tests make the examples trustworthy.

\begin{detailbox}{A Useful Form Review Sentence}
For each field, say: this is the raw editor state, this is the semantic value,
this is the converter, this is the validation policy, this is the problem
target, this is the baseline, and this is how the result uses the value. If the
sentence cannot be spoken, the form is not yet designed.
\end{detailbox}

The migration endpoint is not a form with no Swing code. Swing still owns
documents, focus, editors, keyboard traversal, accessible contexts, and
painting. MigLayout still owns geometry. FlatLaf still paints components.
Corusco owns model facts and binding lifecycles. The mature form is not less
Swing; it is more honest Swing. It lets Swing do widget work and lets the form
model do form work.

This distinction keeps the chapter connected to the rest of the book. The
component chapter explains why text fields, documents, focus, and editors have
their own contracts. The MigLayout chapter explains why the panel should remain
handwritten geometry. The command chapter explains why OK and Save are command
facts. The lifecycle chapter explains why bindings and async validation must
close. The dialog chapter explains the commit protocol. This chapter supplies
the form facts those chapters need.

A book example for this chapter should therefore be more than a pretty dialog.
It should present at least three named states. The first state is clean and
committable, so the reader sees the ordinary form. The second state contains
invalid intermediate raw text, so the reader sees that the text field can
display what the user typed while the semantic value remains protected. The
third state contains a semantic validation problem, so the reader sees that a
parsed value can still be unacceptable. If async validation is demonstrated, a
fourth state should show pending validation and then a stale-result scenario in
a model test. A screenshot without those model stories would be decorative.

The companion code should keep the same split. A core test constructs the form
model directly, edits raw text, invokes validators, and asserts problems. A
Swing smoke test constructs the MigLayout panel, installs bindings in a scope,
types or sets component values on the EDT, and asserts that model state changes.
The screenshot harness fixes size, Look-and-Feel, sample data, and initial
state. These layers do not duplicate work. They prove different claims: the
model owns the facts, the binding connects the facts, and the visual example
communicates the facts.

Documentation should also avoid showing a false shortcut. It is tempting to
quote only the final dialog code because that is what the reader sees. But the
dialog is the least important part if the model is wrong. The chapter's
examples should quote the field model, a converter choice, a validation rule,
and the OK protocol before dwelling on borders or icons. The visual treatment
then feels earned. The reader can trace a red border back to a problem target,
a disabled OK button back to severity, and a dirty marker back to baseline
state.

Operational review is the last pass. Ask what happens if the repository is
slow, if the user closes the dialog during validation, if a translated message
is longer than expected, if a screen reader asks for the field description, if
the user presses Escape with dirty but unparsable text, if Apply succeeds but
OK later fails, and if a save failure returns server-side field problems. None
of these questions belongs solely to Swing or solely to Corusco. They are form
workflow questions. A mature form model gives each answer a place to live.

Server-side validation is the common final check. A local form may be perfectly
committable and still receive a save response saying that a name was taken, a
credit limit exceeded a permission boundary, or an address failed a current
service rule. Treat those responses as problems when the user can repair them.
Map server field identifiers to \api{FieldKey} values, map row failures to row
or cell targets, preserve severity, and keep diagnostic details for logs. Do
not collapse every save failure into a modal "save failed" box. Some failures
belong in a message box; repairable form failures belong back in the form's
problem language so the user can continue editing. That mapping should itself
be covered by tests, because service contracts drift just as surely as Swing
screens do.

\section{Exercises}

\begin{exercisebox}
\begin{enumerate}
\item Take a numeric text field and write down its raw text, semantic value,
parse state, parse problem, dirty state, and touched state after three edits.
\item Replace one error label with a \api{ProblemSet}. Bind the label only as a
presentation of the most severe problem.
\item Write a validation rule that depends on two fields. Record the dependency
explicitly.
\item Add a model-level test for invalid intermediate input.
\item Design an OK flow for a dialog that includes active-editor commit,
validation focus, and dirty cancel confirmation.
\item Add stale-result suppression to an asynchronous uniqueness validator.
\end{enumerate}
\end{exercisebox}

\textbf{Checklist.}

\begin{itemize}
\item Keep raw text, parsed value, validation result, touched state, and dirty
state distinct.
\item Use \api{TextFieldModel} when text edits a semantic value.
\item Use \api{FieldModel} when the component already edits the semantic type.
\item Use stable problem codes and typed targets.
\item Do not throw exceptions for normal user input mistakes.
\item Base committability on problem severity, not label contents.
\item Cancel or ignore stale async validations.
\item Commit active editors before dialog OK validation.
\item Test form models without Swing before testing bindings.
\end{itemize}
